\documentclass[11pt]{article}
\PassOptionsToPackage{english}{babel} % le résumé anglais (otherlanguage) exige la langue english ; le français reste la langue principale
% Préambule commun des prétirages v2 (format arXiv : source TeX autonome, pdfLaTeX).
% Moteur : pdfLaTeX (arXiv) → fontenc T1 + inputenc ; XeTeX/LuaTeX (tectonic) → fontspec. Sans ce choix, les
% guillemets « » tapés en dur sortent en « ń ż » sous tectonic (ledger, 2026-11-14).
\usepackage{iftex}
\ifPDFTeX
  \usepackage[T1]{fontenc}
  \usepackage[utf8]{inputenc}
\else
  \usepackage{fontspec}
\fi
\usepackage[french]{babel}
\usepackage{lmodern,microtype}
\usepackage[a4paper,margin=2.5cm]{geometry}
\usepackage{amsmath,amssymb}
\usepackage{graphicx,booktabs,siunitx}
\usepackage{authblk}
\usepackage[numbers,sort&compress]{natbib}
\usepackage[hidelinks]{hyperref}
\usepackage{csquotes}
\sisetup{locale=FR,group-minimum-digits=4}
\renewcommand\Affilfont{\small}
% Étiquette de provenance obligatoire sur chaque nombre (règle du projet) :
\newcommand{\mesure}{\textsuperscript{\textsc{m}}}   % mesuré
\newcommand{\derive}{\textsuperscript{\textsc{d}}}   % dérivé
\newcommand{\suppose}{\textsuperscript{\textsc{s}}}  % supposé

\usepackage{pgfplots}
\usepgfplotslibrary{fillbetween}
\usepackage{xurl} % coupe les longues URL de la bibliographie
\pgfplotsset{compat=1.16}

\title{Le secteur public était déjà dans les fichiers : séparer public et privé, dater les réformes, coder le zonage et fixer le grain communal d'un jumeau numérique de l'offre de soins}
\author[1]{Équipe hackathon-santé}
\affil[1]{Projet « jumeau numérique de l'offre de soins », dépôt \texttt{hackathon-sante/v2}, axe~4 (données)}
\date{Prétirage, version du 14 novembre 2026 --- non relu par des pairs}

\begin{document}
\maketitle

\begin{abstract}
Le jumeau numérique de l'offre de soins projette jusqu'en 2050 l'effectif « libéral » de cinq professions
(sages-femmes, chirurgiens-dentistes, pharmaciens, masseurs-kinésithérapeutes, infirmiers). Ici, « libéral »
désigne le secteur privé, lucratif ou non : tout employeur public est exclu. Or le stock DREES utilisé jusqu'ici
ne séparait pas public et privé chez les salariés. La sortie 2026 n'était donc connue qu'en fourchette, de
\num{12475} à \num{26428} sages-femmes. Nous montrons que la séparation est déjà dans les fichiers ouverts de la
DREES : ils croisent département, secteur d'activité, mode d'exercice, âge et sexe. Les salariés des secteurs
mixtes sont répartis à l'aide du répertoire FINESS. La fourchette des sages-femmes devient
\num{12059}--\num{14660}\mesure{} : sa largeur baisse de 81\,\%\derive{}. Les infirmiers sont séparés sans
ambiguïté : \num{277788} privés en 2023\mesure{}. Nous livrons aussi un registre de 30 événements datés et
sourcés, et le zonage conventionnel de \num{34919} géocodes pour quatre professions. Chaque table est contrôlée
contre un total connu. Le registre est confirmé sur un cas : l'année sans diplômés en kinésithérapie (2018) se
lit dans le stock, avec \num{-35}\,\% de moins de 25 ans\mesure{}. Le zonage reproduit le périmètre dentaire
publié, soit 97 géocodes et 4,8\,\% de la population\mesure{}. Enfin, l'Annuaire Santé rend les cinq
professions observables à la commune. Mais la médiane des effectifs privés par commune est nulle. Le grain
projetable est donc le bassin de vie ou l'EPCI pour les infirmiers, et le département pour les dentistes et les
sages-femmes. Ce même fichier montre que 24\,\% des dentistes et 32\,\% des kinésithérapeutes actifs n'ont pas
de diplôme d'État français\mesure{}. Limite : les infirmiers n'entrent dans le jumeau que sous l'hypothèse d'un
secteur indépendant de l'âge et du sexe, testée sur le sexe seulement.
\end{abstract}
\begin{otherlanguage}{english}
\begin{abstract}
The digital twin of French healthcare supply projects, to 2050, the ``liberal'' headcount of five professions:
midwives, dentists, pharmacists, physiotherapists and nurses. Here ``liberal'' means the private sector, for-profit
or not, with every public employer excluded. The DREES stock used so far did not separate public from private
employers among salaried staff, so the 2026 output was only known as a range: 12,475 to 26,428 midwives. We show
that the split is already present in the open DREES files, which cross département, activity sector, practice
mode, age and sex. Salaried staff in mixed sectors are apportioned using the FINESS register. The midwife range
narrows to 12,059--14,660 (measured), 81\,\% narrower. Nurses separate without ambiguity: 277,788 in the private
sector in 2023. We also release a register of 30 dated, sourced events and the conventional zoning of 34,919
geocodes for four professions. Each table is checked against a known total. One event is confirmed in the
data: the 2018 year without physiotherapy graduates appears in the stock, with 35\,\% fewer under-25s. The
zoning reproduces the published dental perimeter: 97 geocodes and 4.8\,\% of the population. Finally, the public
Annuaire Santé register makes all five professions observable by municipality. The median private headcount per
municipality is zero, however. The projectable grain is therefore the living area or inter-municipal body
for nurses, and the département for dentists and midwives. The same register shows that 24\,\% of active dentists
and 32\,\% of active physiotherapists hold no French state diploma. Limitation: nurses enter the twin only under
the assumption that sector is independent of age and sex, which is tested on sex only.
\end{abstract}
\end{otherlanguage}
\noindent\textbf{Mots-clés :} démographie des professions de santé ; secteur public et privé ; RPPS ; FINESS ;
zonage conventionnel ; régulation de l'installation ; données ouvertes ; jumeau numérique.

\noindent\textbf{Provenance des nombres :} \mesure{}~mesuré sur les données, \derive{}~dérivé par un calcul
explicite, \suppose{}~supposé.

\section{Introduction}
Le projet simule des politiques nationales sur un jumeau de l'offre de soins. Chaque politique y est un
réglage d'hyperparamètres sourcés. La sortie est l'effectif « libéral » de cinq professions, à toute échelle
jusqu'à la commune. Il s'agit de la définition du projet, qui diffère de celle des sources : le secteur privé,
lucratif ou non, y compris les pharmaciens adjoints d'officine, les salariés des cliniques et l'industrie. Tout
employeur public en est exclu : hôpital public, PMI, Éducation nationale, EHPAD publics, centres municipaux.

Trois lacunes de données bloquaient ce cadrage :
\begin{enumerate}
  \item le stock ne séparait pas public et privé dans « salarié hospitalier » et « autre salarié ». La sortie
  2026 n'était donc connue qu'en fourchette : \num{12475}--\num{26428} sages-femmes et
  \num{65344}--\num{75335} pharmaciens\mesure{} ;
  \item les ruptures des séries (formation, diplômés, installation) n'étaient rattachées à aucune réforme datée ;
  \item la régulation de l'installation, qui s'applique par commune, n'existait dans le jumeau que sous forme de
  zones illustratives.
\end{enumerate}

\paragraph{Contributions.}
\begin{itemize}
  \item Une séparation public / privé par profession, mode, département, âge, sexe et année (2012--2026). Elle
  donne pour chaque sortie une fourchette resserrée et une valeur centrale (section~\ref{sec:res-secteur}).
  \item Un registre de 30 événements datés, chacun avec sa source ouverte. Il comprend deux ruptures de données,
  dont un artefact de fichier qui passait pour une réforme (section~\ref{sec:res-evenements}).
  \item Le zonage conventionnel de \num{34919} géocodes pour quatre professions, sur deux millésimes, avec le
  statut lu par le module de régulation (section~\ref{sec:res-zonage}).
  \item Un inventaire des sources ouvertes à la commune. Il comble quatre trous : Marseille, sages-femmes,
  pharmaciens et infirmiers par âge et sexe. Il donne aussi le grain projetable par profession
  (section~\ref{sec:res-grain}).
\end{itemize}

\section{Travaux existants}
La DREES publie chaque année les effectifs au 1\textsuperscript{er} janvier, à partir du RPPS
\citep{drees2026demo}. La feuille « Lisez-moi » documente le croisement par département, secteur d'activité,
mode d'exercice, âge et sexe. Les kinésithérapeutes y sont comptés au RPPS depuis 2017, et à partir d'ADELI
avant cette date \citep{drees2020er1159}. Pour les infirmiers, la DREES croise deux sources : la Base Tous Salariés
(Insee) pour les salariés, le SNDS pour les libéraux \citep{dreesbts2026}. Le répertoire FINESS décrit les
établissements et le statut juridique de leurs gestionnaires \citep{finessEt2026,finessEj2026}.

Le zonage conventionnel est arrêté par les ARS et diffusé par CartoSanté \citep{cartosante2026}. Une génération
antérieure est publiée par le catalogue Atlasanté \citep{atlasante2014zonage}. Les règles qui s'y attachent
viennent des avenants aux conventions nationales :
\begin{itemize}
  \item infirmiers \citep{ameliIde} ;
  \item sages-femmes \citep{ameliSf} ;
  \item kinésithérapeutes \citep{ameliMk}, dont la régulation a été annulée en 2014 par le Conseil d'État
  \citep{ce2014mk} ;
  \item chirurgiens-dentistes \citep{arspaca2025cd}.
\end{itemize}
L'Agence du numérique en santé publie chaque jour une extraction en libre accès du RPPS \citep{ans2026annuaire}.
Elle donne la commune de chaque activité, mais ni l'âge ni d'archive. L'Insee publie la composition communale
des bassins de vie 2022 \citep{insee2026bv}.
Nous n'avons pas trouvé de travail qui exploite le croisement secteur~$\times$~mode du fichier DREES pour isoler
le secteur privé au sens large.

\section{Données}
\begin{table}[h]
\centering\small
\begin{tabular}{@{}p{4.2cm}p{3.1cm}p{4.3cm}p{3.2cm}@{}}
\toprule
Jeu & Producteur, millésime & Grain & Usage \\
\midrule
Démographie des PS depuis 2012 \citep{drees2026demo} & DREES, 2026-07-02 & dép.\ $\times$ secteur (9--12) $\times$ mode (4) $\times$ âge (10) $\times$ sexe, 2012--2026 & SF, CD, PH, MK \\
Infirmières salariées BTS, libérales SNDS \citep{dreesbts2026} & DREES / Insee, 2026-07-02 & dép.\ $\times$ secteur (6), sans croisement âge ou sexe ; 2013--2023 & IDE \\
FINESS établissements et entités juridiques \citep{finessEt2026,finessEj2026} & Ministère, au 2026-05-04 & établissement, catégorie, statut juridique & parts privées \\
CartoSanté \texttt{zonage\_conv} \citep{cartosante2026} & CNAMTS--ARS, extrait 2026-11-14 & commune COG 2026, arrondissements PLM & zonage actuel \\
Atlasanté c987d0e2 \citep{atlasante2014zonage} & ARS, fiche 2014 & commune (\num{36700}) & zonage 2012 \\
Découpage administratif \citep{geoapi2026} & Etalab, 2026-11-14 & commune, population légale & contrôles \\
\bottomrule
\end{tabular}
\caption{Jeux de données. Tous sont ouverts. Les fichiers FINESS pèsent \SI{35.8}{\mega\byte} et
\SI{11.9}{\mega\byte}\mesure{}, le zonage environ \SI{1.6}{\mega\byte} par profession\mesure{}.}
\label{tab:donnees}
\end{table}

\paragraph{Contrôles d'entrée.}
\begin{itemize}
  \item La somme des cellules détaillées (dép.~$\times$ secteur $\times$ mode $\times$ âge $\times$ sexe) égale
  le total France du fichier, à 0\,\% près\mesure{} pour les sages-femmes, les dentistes, les pharmaciens et
  ADELI.
  \item Les kinésithérapeutes RPPS font exception : 0,22\,\% d'écart au plus\mesure{}, en 2017--2018. L'écart est
  interne au fichier source : en 2017, la somme des secteurs vaut \num{84599} contre \num{84521} pour
  « Ensemble »\mesure{}.
  \item Pour le BTS, l'écart est d'au plus 0,009\,\%\mesure{}.
\end{itemize}

\section{Méthode}
\subsection{Classement public / privé}
Soit $N_{p,s,m,d,a,x,t}$ l'effectif de la profession $p$ dans le secteur $s$, le mode $m$, le département $d$,
l'âge $a$, le sexe $x$ et l'année $t$. Chaque cellule reçoit une part privée $\pi \in [0,1]$ :
\begin{itemize}
  \item libéral ou mixte : $\pi = 1$, quel que soit le secteur. C'est la définition du projet, dans laquelle
  les mixtes sont inclus.
  \item Salarié dont le secteur est sans ambiguïté : $\pi = 0$ pour l'hôpital public ; $\pi = 1$ pour les
  établissements ESPIC et privés lucratifs, l'officine, les laboratoires, le médicament et les cabinets.
  \item Salarié d'un secteur mixte (centre de santé, prévention, médico-social, « autre ») : $\pi$ est inconnu.
  Il est borné par $[0,1]$ et prend une valeur centrale $\hat\pi$.
\end{itemize}
La sortie privée est alors encadrée et centrée :
\begin{equation}
N^{\min}_{p,t} = \sum_{\text{sûr privé}} N, \qquad
N^{\max}_{p,t} = N^{\min}_{p,t} + \sum_{\text{ambigu}} N, \qquad
\hat N_{p,t} = \sum \hat\pi\, N .
\end{equation}

\subsection{Parts centrales par FINESS}
La valeur $\hat\pi$ est la part d'établissements de la catégorie FINESS correspondante dont l'entité juridique est
privée :
\begin{itemize}
  \item statuts juridiques 47 à 99 : privé (mutuelles, associations, fondations, sociétés) ;
  \item statuts 1 à 39 : public ;
  \item statuts 40 à 46, organismes de sécurité sociale : comptés publics, comme administrations de sécurité
  sociale.
\end{itemize}
Les catégories retenues sont :
\begin{itemize}
  \item centres de santé (124) ; pour les dentistes, les centres dont la raison sociale contient « DENT » ;
  \item PMI et planification (223, 228) ;
  \item EHPAD (500) ;
  \item handicap médicalisé (182, 183, 186, 190, 255, 437, 448).
\end{itemize}
« Autre secteur », et « prévention et établissement de santé » chez les pharmaciens, n'ont pas de proxy :
$\hat\pi = 0{,}5$\suppose{}.

\paragraph{Infirmiers.} Le BTS sépare déjà public et privé en trois paires : hôpital, médico-social, autre. Les
libéraux SNDS sont privés. Aucune ambiguïté n'est donc introduite. Le BTS ne croise pas le secteur avec l'âge ou
le sexe, et la table infirmiers n'a que le grain département~$\times$ secteur.

\subsection{Registre d'événements}
Chaque ligne de \texttt{v2/data/events.csv} est un texte ou un fait observé. Elle porte ces champs :
\texttt{id}, \texttt{date\_effet}, \texttt{profession}, \texttt{flux\_touche} (capacité de formation, durée des
études, diplômés, installation, régulation, compétences, autre), \texttt{type}, \texttt{description},
\texttt{source}, \texttt{millesime} et \texttt{confiance}.
\begin{itemize}
  \item Chaque URL a été ouverte, et le fait lu sur la page.
  \item \texttt{verifie} : date et fait lus sur la source. \texttt{partiel} : une partie de la date ou du chiffre
  est déduite, et la description le dit.
  \item Quand seule l'année est connue, la date retenue est la rentrée (1\textsuperscript{er} septembre) ou la fin
  d'année universitaire.
\end{itemize}

\subsection{Zonage et statut}
Le statut lu par le module de régulation se déduit de la classe ARS :
\begin{itemize}
  \item Statut \texttt{remplacement\_seul} : classes « sur-dotée » (ancien zonage) et « non prioritaire »
  (nouveau zonage). Le conventionnement n'y est possible qu'en remplacement d'un départ.
  \item Très sous-dotée ou sous-dotée : \texttt{incitation}.
  \item Les autres classes : \texttt{autre}.
\end{itemize}
Le référentiel \texttt{arm2026} remplace Paris, Lyon et Marseille par leurs 45 arrondissements, zonés à ce niveau.

\section{Résultats}
\subsection{Séparation public / privé}\label{sec:res-secteur}
\begin{table}[h]
\centering\small
\setlength{\tabcolsep}{4pt}
\begin{tabular}{@{}lrrrrrr@{}}
\toprule
Profession & Année & Effectif & Ancienne & Min & Centrale & Max \\
 & & & fourchette & & & \\
\midrule
Sages-femmes & 2026 & \num{26428} & \num{12475}--\num{26428} & \num{12059} & \num{12912} & \num{14660} \\
Pharmaciens  & 2026 & \num{75335} & \num{65344}--\num{75335} & \num{60005} & \num{63612} & \num{67219} \\
Dentistes    & 2026 & \num{48690} & \num{48143}--\num{48690} & \num{40784} & \num{47349} & \num{47894} \\
Kinés        & 2026 & \num{111141} & \num{98348}--\num{111237} & \num{98828} & \num{101951} & \num{103986} \\
Infirmiers   & 2023 & \num{607644} & --- & \num{277788} & \num{277788} & \num{277788} \\
\bottomrule
\end{tabular}
\caption{Effectif « libéral » (= privé), France entière. Min, max et effectif sont mesurés\mesure{} ; la centrale
dépend des parts FINESS (mesurées\mesure{}) et de trois parts supposées à 0,5\suppose{}. Pour les infirmiers, le
total additionne BTS et SNDS (double compte possible, voir section~\ref{sec:menaces}).}
\label{tab:fourchettes}
\end{table}

La largeur de la fourchette baisse de 81\,\% pour les sages-femmes (de \num{13953} à \num{2601}), de 60\,\% pour
les kinésithérapeutes et de 28\,\% pour les pharmaciens\derive{}.

\paragraph{Les anciennes bornes basses n'étaient pas des minima.} Elles comptaient tout « autre salarié » comme
privé, PMI et centres municipaux compris. C'est pourquoi le minimum baisse pour les sages-femmes et les
pharmaciens. Pour les dentistes, la fourchette s'élargit. Les \num{6694} dentistes salariés des centres de santé
sont désormais ambigus\mesure{}, mais leur part centrale est de 0,963\mesure{} : 96\,\% des \num{1299} centres
« dentaires » ont un gestionnaire privé.

\begin{table}[h]
\centering\footnotesize
\begin{tabular}{@{}llrrl@{}}
\toprule
Profession & Secteur & Effectif 2026 & $\hat\pi$ & Source de $\hat\pi$ \\
\midrule
Pharmaciens & Autre secteur & \num{5805} & 0,50\suppose & --- \\
Pharmaciens & Prévention et étab.\ de santé & \num{1409} & 0,50\suppose & --- \\
Dentistes & Centre de santé & \num{6694} & 0,963\mesure & FINESS, \num{1299} centres \\
Kinés & Étab.\ sanitaire et social & \num{2200} & 0,696\mesure & FINESS, \num{13407} étab. \\
Kinés & Autre secteur & \num{2524} & 0,50\suppose & --- \\
Sages-femmes & Prévention (PMI\ldots) & \num{1329} & 0,070\mesure & FINESS, \num{1435} étab. \\
Sages-femmes & Autre secteur & \num{796} & 0,50\suppose & --- \\
\bottomrule
\end{tabular}
\caption{Origine de l'ambiguïté restante : salariés des secteurs mixtes, 2026 (lignes de plus de 500 personnes).}
\label{tab:ambigu}
\end{table}

\begin{figure}[h]
\centering
\begin{tikzpicture}
\begin{axis}[width=0.85\textwidth,height=6cm,xlabel={Année (au 1\textsuperscript{er} janvier)},
  ylabel={Sages-femmes « libérales » (privé)},xmin=2012,xmax=2026,
  x tick label style={/pgf/number format/1000 sep=},
  y tick label style={/pgf/number format/1000 sep={\,}},
  legend pos=north west,legend style={font=\small}]
\addplot[name path=max,draw=none,forget plot] coordinates {(2012,8337) (2013,9015) (2014,9568) (2015,10184) (2016,10864) (2017,11296) (2018,11614) (2019,11898) (2020,12254) (2021,12692) (2022,13047) (2023,13400) (2024,13878) (2025,14214) (2026,14660)};
\addplot[name path=min,draw=none,forget plot] coordinates {(2012,6839) (2013,7675) (2014,8127) (2015,8704) (2016,9338) (2017,9760) (2018,10024) (2019,10255) (2020,10524) (2021,10707) (2022,10962) (2023,11238) (2024,11510) (2025,11706) (2026,12059)};
\addplot[gray!30] fill between[of=max and min];
\addlegendentry{min--max (secteurs mixtes)}
\addplot[thick,black,mark=*,mark size=1.2pt] coordinates {(2012,7309) (2013,7980) (2014,8459) (2015,9020) (2016,9656) (2017,10075) (2018,10356) (2019,10607) (2020,10929) (2021,11239) (2022,11527) (2023,11857) (2024,12233) (2025,12500) (2026,12912)};
\addlegendentry{centrale (FINESS)}
\end{axis}
\end{tikzpicture}
\caption{Sages-femmes du secteur privé, France entière, 2012--2026\mesure{}. L'ancienne fourchette 2026
(\num{12475}--\num{26428}) dépasserait le haut du graphique.}
\label{fig:sf}
\end{figure}

Les séries existent pour chaque année et chaque département. La largeur médiane de la fourchette par département,
dernière année, va de 4\,\% de l'effectif (kinésithérapeutes) à 9\,\% (sages-femmes, dentistes, pharmaciens)\mesure{}.

\subsection{Registre d'événements}\label{sec:res-evenements}
Le registre compte 30 événements entre 2009 et 2028\mesure{} : 23 \texttt{verifie} et 7 \texttt{partiel}.
\begin{table}[h]
\centering\small
\begin{tabular}{@{}lrp{10.5cm}@{}}
\toprule
Flux touché & N & Exemples (date d'effet) \\
\midrule
Régulation de l'installation & 11 & IDE 18/04/2009, 27/05/2012 \citep{ameliIde} ; MK 14/01/2012, annulée le 17/03/2014 \citep{ce2014mk}, puis 01/07/2018 et avenant~7 (2023) \citep{ameliMk} ; SF 15/09/2012 \citep{ameliSf} ; CD 01/01/2025 \citep{arspaca2025cd} \\
Capacité de formation & 7 & numerus apertus, rentrée 2020 \citep{loiotss2019} ; objectifs pluriannuels 2021--2025 et 2026--2030 ; 8 sites d'odontologie à la rentrée 2022 \\
Durée des études & 4 & MK : 3 à 4 ans en IFMK, rentrée 2015 \citep{arretemk2015} ; SF : 6 ans, entrée en 2\textsuperscript{e} année à partir de 2024 \\
Diplômés & 2 & MK : année blanche 2018 \citep{drees2020er1159} ; SF : (quasi) aucune diplômée en 2027-28 (déduit) \\
Compétences & 3 & IPA (2018) ; vaccination en officine (2019, 2023) \\
Installation (officines) & 1 & seuils de population (ordonnance 2018-3) \\
Données & 2 & passage ADELI $\rightarrow$ RPPS des MK en 2017 \citep{drees2020er1159} ; artefact des inscrits 2025-26 \\
\bottomrule
\end{tabular}
\caption{Événements par flux touché. Liste complète, sources et fiches : \texttt{v2/data/events.csv},
\texttt{v2/docs/lab/fiches/data-ev-*.md}.}
\label{tab:evenements}
\end{table}

\paragraph{Le saut d'inscrits 2025 est un artefact.} Le fichier des inscrits 2025-26 compte \num{2083} étudiants
en 2\textsuperscript{e} année d'odontologie\mesure{}, contre environ \num{1550} les années précédentes. Une seule
ligne explique le saut : l'Université Paris Cité, avec 683 inscrits\mesure{}, contre 188 à 215 entre 2020 et 2024
et 193 en 3\textsuperscript{e} année. Le fichier compilé de l'équipe donne \num{1588} au national, dans la tendance
(\num{1546} en 2024). Le registre enregistre donc ce saut comme une rupture de données, pas comme une réforme.

\subsection{Zonage par commune}\label{sec:res-zonage}
\begin{table}[h]
\centering\small
\begin{tabular}{@{}lrrrr@{}}
\toprule
Profession & Remplacement seul & Incitation & Autre & Non renseigné \\
\midrule
Sages-femmes & \num{1471} \;(12,3\,\%) & \num{11862} \;(17,4\,\%) & \num{21562} \;(70,3\,\%) & 24 \\
Dentistes & 97 \;(4,8\,\%) & \num{26769} \;(44,4\,\%) & \num{8027} \;(50,8\,\%) & 26 \\
Kinés & \num{2836} \;(27,4\,\%) & \num{18840} \;(29,8\,\%) & \num{13241} \;(42,8\,\%) & 2 \\
Infirmiers & \num{4786} \;(27,6\,\%) & \num{5690} \;(15,3\,\%) & \num{24441} \;(57,0\,\%) & 2 \\
\bottomrule
\end{tabular}
\caption{Zonage actuel (CartoSanté, extrait le 14/11/2026) : nombre de géocodes et part de la population
légale\mesure{}. Les géocodes non renseignés sont des communes des DROM. Ils sont comptés avec « autre » dans les
parts de population.}
\label{tab:zonage}
\end{table}

\begin{itemize}
  \item Chez les kinésithérapeutes, deux régimes coexistent : \num{2826} géocodes « non prioritaires » (nouveau
  zonage) et 10 « sur-dotés » (ancien)\mesure{}.
  \item La génération 2012 (métropole, sans sages-femmes) classe en remplacement seul \num{3135} communes pour
  les infirmiers, 547 pour les kinésithérapeutes et 908 pour les dentistes\mesure{}.
\end{itemize}

\subsection{Grain communal}\label{sec:res-grain}
\paragraph{Inventaire.} Trois sources ouvertes descendent à la commune.
\begin{itemize}
  \item L'Annuaire Santé couvre les cinq professions : un instantané du 13 novembre 2026, avec la civilité, le mode,
  le secteur, le FINESS et l'origine du diplôme, mais sans âge \citep{ans2026annuaire}.
  \item CartoSanté donne l'histoire des libéraux par tranche d'âge : infirmiers 2004--2025, kinésithérapeutes
  2005--2025, dentistes 2006--2025 \citep{cartosante2026}.
  \item FINESS décrit les structures.
\end{itemize}
Deux pistes ont été écartées :
\begin{itemize}
  \item les séries « flux » de CartoSanté ne sont pas des installations, mais des pôles d'attraction des patients ;
  \item la série des sages-femmes y est illisible : les modalités d'âge sont libellées en années.
\end{itemize}

\paragraph{Trous comblés.}
\begin{itemize}
  \item \textbf{Marseille.} Ses 16 arrondissements couvrent 2010--2025. Pour les infirmiers, leur somme reproduit
  l'écart du département 13 à l'unité près, chaque année\mesure{}.
  \item \textbf{Pharmaciens et sages-femmes.} L'Annuaire donne leur instantané communal : \num{61882}
  pharmaciens et \num{11612} sages-femmes privés\mesure{}. Le premier chiffre tombe dans la fourchette DREES
  2026 (table~\ref{tab:fourchettes}).
  \item \textbf{Infirmiers.} Ils sont ventilés par secteur, âge et sexe sous l'hypothèse $H_{\mathrm{ind}}$
  (secteur indépendant de l'âge et du sexe dans un département). Le test, fixé avant de voir le résultat,
  compare la part de femmes du privé et du public. L'écart médian est de 2,1 points par département, sous le
  seuil de 3 points : l'hypothèse n'est pas rejetée\mesure{}.
\end{itemize}

\paragraph{Contrôles.} La somme des communes CartoSanté est rapprochée du RPPS libéral et mixte de la même année.
En 2025, l'écart est de \num{-3.2}\,\% pour les kinésithérapeutes et de \num{-0.8}\,\% pour les dentistes\mesure{}.
Pour les infirmiers, l'écart au SNDS est de +4,7\,\% en 2023\mesure{}.

\begin{table}[h]
\centering\small
\begin{tabular}{@{}lrrrrr@{}}
\toprule
Profession & Commune & Bassin de vie & EPCI & Département & Pros dans une commune $<30$ \\
\midrule
Infirmiers & 0 & 35,5 & 53 & \num{1644} & 30,6\,\% \\
Kinés & 0 & 17 & 23 & 615 & 49,8\,\% \\
Pharmaciens & 0 & 11 & 16,5 & 444 & 57,5\,\% \\
Dentistes & 0 & 6 & 8 & 261 & 53,4\,\% \\
Sages-femmes & 0 & 1,5 & 2,5 & 78 & 69,8\,\% \\
\bottomrule
\end{tabular}
\caption{Médiane des effectifs privés par unité (Annuaire Santé, 13/11/2026)\mesure{}. Unités : \num{34917}
communes (arrondissements PLM), \num{1707} bassins de vie \citep{insee2026bv}, \num{1255} EPCI
\citep{geoapi2026}, 101 départements.}
\label{tab:grain}
\end{table}

\paragraph{Grain projetable.} La règle retenue est conventionnelle : une médiane d'au moins 30 professionnels
privés par unité, soit un bruit de Poisson d'au plus 18\,\%\derive{}. Elle donne :
\begin{itemize}
  \item infirmiers : bassin de vie ou EPCI ;
  \item kinésithérapeutes et pharmaciens : EPCI, avec un modèle hiérarchique qui emprunte de la force au
  département ;
  \item dentistes et sages-femmes : département.
\end{itemize}
La commune reste un grain d'observation et de répartition. Aucune source ouverte ne donne de flux d'installation
par commune. Le meilleur substitut est la variation du stock des moins de 40 ans dans CartoSanté.

\paragraph{Diplômés hors de France.} Part des actifs sans diplôme d'État français\mesure{} : 24,4\,\% des
dentistes et 32,5\,\% des kinésithérapeutes, contre 9,7\,\% (sages-femmes), 3,5\,\% (pharmaciens) et
2,6\,\% (infirmiers). Ce sont surtout des diplômes de l'UE.

\section{Validation et menaces à la validité}\label{sec:menaces}
\paragraph{Contrôles contre des totaux connus.}
\begin{itemize}
  \item \emph{Stock} : la somme des cellules égale le total publié, avec un écart d'au plus 0,22\,\%\mesure{}.
  Cet écart est interne à la source.
  \item \emph{Événements} : l'année blanche des kinésithérapeutes (2018) doit se voir au 1\textsuperscript{er}
  janvier 2019. C'est le cas : les moins de 25 ans passent de \num{5637} à \num{3681} (\num{-35}\,\%), puis
  remontent à \num{5126} en 2026\mesure{}.
  \item \emph{Zonage} :
  \begin{itemize}
    \item communes : le COG en compte \num{34969} ; moins 92 communes des COM hors champ, moins les 3 communes
    PLM, plus 45 arrondissements, on obtient \num{34919} géocodes, comme dans la table\mesure{} ;
    \item périmètre dentaire : la table donne 97 géocodes et 4,8\,\% de la population\mesure{}. Les sources
    publiées citent 5\,\% de la population \citep{arspaca2025cd} et 97 communes\footnote{Mémoire EHESP 2025,
    § 3.3.2, déposé dans le dépôt du projet et lu par la session de l'axe~1.}.
  \end{itemize}
\end{itemize}

\paragraph{Ce qui pourrait rendre les résultats faux.}
\begin{itemize}
  \item \textbf{Secteur de l'activité la plus récente.} La DREES classe chaque professionnel selon son activité la
  plus récente, et non selon son activité principale \citep{drees2026demo}. Un salarié public qui prend un second
  poste privé est donc compté privé.
  \item \textbf{Parts en établissements.} Les parts FINESS comptent des établissements, pas des salariés. Un grand
  EHPAD public pèse autant qu'un petit EHPAD privé. L'écart n'est pas mesuré.
  \item \textbf{Trois parts supposées.} Elles portent l'essentiel de l'incertitude des pharmaciens : \num{7214}
  salariés de secteurs mixtes, soit 9,6\,\% de l'effectif\derive{}.
  \item \textbf{Recoupement avec l'Ordre.} Parmi les pharmaciens hospitaliers, la part privée vaut 17\,\% selon la
  DREES et 28,6\,\% selon l'Ordre (inscriptions en section H). Ce chiffre est repris des fiches de l'axe~1, et
  l'écart n'est pas réconcilié.
  \item \textbf{Infirmiers.}
  \begin{itemize}
    \item Un double compte est possible entre le SNDS et le BTS.
    \item Le BTS a une rupture de série en 2022.
    \item \emph{Surtout, le BTS ne croise pas le secteur avec l'âge ni le sexe.} Les infirmiers n'entrent
    donc dans le jumeau que sous $H_{\mathrm{ind}}$. L'hypothèse a été testée sur le sexe (non rejetée), pas sur
    l'âge : des jeunes plus souvent à l'hôpital public la fausseraient.
  \end{itemize}
  \item \textbf{Zonage.}
  \begin{itemize}
    \item Le millésime actuel n'est pas daté : CartoSanté ne sert pas la date des arrêtés.
    \item L'historique n'a qu'une génération (2012).
    \item L'« obligation d'exercice » prévue pour les kinésithérapeutes (avenant~7) n'est pas une classe de zonage.
  \end{itemize}
  \item \textbf{Grain communal.}
  \begin{itemize}
    \item L'Annuaire est un instantané de novembre, la DREES compte au 1\textsuperscript{er} janvier. 11,5\,\% des
    activités dentaires n'y ont pas de commune.
    \item Le seuil de 30 est une convention, que seul le banc peut remplacer.
  \end{itemize}
  \item \textbf{Registre.}
  \begin{itemize}
    \item 7 lignes sont partielles.
    \item L'année sans diplômées SF 2027-28 est déduite du calendrier des études.
    \item L'avenant~7 des kinésithérapeutes prend effet à des dates différentes selon les ARS.
  \end{itemize}
\end{itemize}

\section{Reproductibilité}
Depuis la racine du dépôt, avec l'environnement \texttt{uv} du projet (Python~3.13, pandas, pyarrow,
python-calamine) :
\begin{verbatim}
uv run python v2/data/build_secteur.py   # ~12 s
uv run python v2/data/build_zonage.py    # ~2 s
uv run python v2/data/build_annuaire.py  # ~15 s
uv run python v2/data/build_commune.py   # ~3 s
uv run python v2/data/build_infirmiers.py # ~2 s
\end{verbatim}
\begin{itemize}
  \item \textbf{Entrées.}
  \begin{itemize}
    \item Fichiers DREES dans \texttt{data/raw/drees/} (publication du 2 juillet 2026).
    \item FINESS dans \texttt{v2/data/raw/finess/} (\texttt{etalab-cs1100502} et \texttt{-cs1100501}, au
    2026-05-04).
    \item Zonage dans \texttt{v2/data/raw/zonage/} : réponses TJS \texttt{GetData} du 14 novembre 2026 et GeoJSON
    Atlasanté.
  \end{itemize}
  \item \textbf{Sorties} (\texttt{v2/data/processed/}) :
  \begin{itemize}
    \item \texttt{secteur\_stock.parquet} (\num{899250} lignes) ;
    \item \texttt{secteur\_split.parquet} ;
    \item \texttt{secteur\_checks.csv} et \texttt{secteur\_finess\_shares.csv} ;
    \item \texttt{zonage\_commune.parquet}.
  \end{itemize}
  \item Le registre \texttt{v2/data/events.csv} est versionné tel quel.
  \item Aucun tirage aléatoire n'intervient.
  \item Les rapports détaillés sont dans \texttt{v2/docs/lab/reports/2026-11-14-axe4-*.md}.
\end{itemize}

\section{Conclusion}
\paragraph{Ce que ça change pour les autres axes.}
\begin{itemize}
  \item \textbf{Jumeau (axe~0).} Les données couvrent les deux cas suivants.
  \begin{itemize}
    \item L'état peut porter le secteur public observé (2012--2026) pour quatre professions, au lieu du
    substitut « libéral + mixte ».
    \item Le module de régulation lit des géocodes réels.
  \end{itemize}
  \item \textbf{Banc (axe~3).} Le banc se relance sur la valeur centrale, avec min et max en bande. L'expérience
  naturelle sur le zonage infirmier dispose de la génération 2012 au même format.
  \item \textbf{Filière de formation.} L'artefact 2025 doit être corrigé avant toute projection 2030+. Le flux
  des diplômés hors de France a désormais un ordre de grandeur.
  \item \textbf{Axe~2.} Le grain du modèle hiérarchique est fixé par profession.
\end{itemize}

\paragraph{Questions ouvertes, avec leur coût} (registre \texttt{v2/docs/lab/open.md}) :
\begin{itemize}
  \item ventiler l'« autre secteur » des pharmaciens à l'aide des sections de l'Ordre (\(\approx\)1~h) ;
  \item pondérer les parts FINESS par les effectifs salariés (\(\approx\)1~h~30) ;
  \item archiver l'Annuaire chaque mois, seule voie vers des flux d'installation individuels (\(\approx\)30~min,
  puis temps du monde) ;
  \item dater le zonage par région (\(\approx\)1~h) ;
  \item ouvrir les textes primaires des 12 événements à source secondaire (\(\approx\)45~min).
\end{itemize}

\bibliographystyle{plainnat}
\bibliography{refs}
\end{document}
